%% kapa: exam=2 round=31 year=2020 period=2 subject=appraisal_theory subject_ko=감정평가이론 seq=1 number=1 points=40
%% source: Q-Net 감정평가사 기출문제 (cst00309) artlSeq=5207673 file=raw/1008/5207673/2교시_감정평가이론.pdf
%% conversion: kordoc 4.24.51 -> tools/gen2.py (auto); review=proofread-gpt-astra; flags=none
\begin{longquestion}{【문제 1】}
감정평가와 관련한 다음의 물음에 답하시오. (40점) \par
\medskip\noindent 물음 1) 감정평가의 개념을 구체적으로 설명하고, 감정평가의 개념에 근거하여 기준 가치 확정과 복수(複數) 감정평가의 필요성에 관하여 각각 논하시오. (20점) \par
\medskip\noindent 물음 2) 시장가치와 시장가격(거래가격)의 개념을 비교하여 설명하고, 다양한 제도를 통해 시장가격(거래가격)을 수집ㆍ분석할 수 있음에도 불구하고 감정평가가 필요한 이유에 관하여 논하시오. (20점) \par
\end{longquestion}
